% Data --- Why Banks Open Branches in the Digital Era
% Last edited: 2026-05-12

\section{Data}
\label{sec:data}

\subsection{Data Sources and Sample Construction}
\label{sec:data:sources}

The empirical question asks where, and on what basis, large banks still open de novo branches in a period of net network contraction (Figure~\ref{fig:us_branches_total}). Answering it requires four pieces of measurement: a clean count of de novo openings, a candidate-market set in which non-entry carries information, bank--market proxies for the three channels the paper hypothesizes pull a bank into a market (deposit contestability, cost-side edge, existing non-branch lending), and a vector of local controls that absorbs demand variation.

We build the panel at the bank--ZIP--year level over 2012--2026. Branch counts and bank size come from the FDIC Summary of Deposits; an opening is a bank--ZIP--year in which the bank operates at least one branch whose UNINUMBR is new to the SOD in $t$. The bank universe is restricted to commercial banks and thrifts whose maximum reported assets ever reach one billion dollars; the community-bank tail is excluded because the branching margin is qualitatively different and because Call Report items used downstream are not reliably reported below that threshold. We further exclude RSSDIDs that first appear in the SOD in 2011, when the OCC's absorption of the OTS reset the identifier space and a post-merger thrift is mechanically indistinguishable from a de novo entry. Within each bank--year, the candidate ZIP set is built from the bank's revealed geography. For each county in which the bank operates at least one branch, we include every ZIP in that county's CBSA if the county is part of a CBSA, and every ZIP in the county itself if the county lies outside any CBSA. The bank's candidate ZIP set is the union of these per-operating-county expansions. Approximately eighty-five to ninety percent of the de novo openings observed in the SOD over the sample period fall inside this set, so restricting the opportunity set to it preserves the variation the entry margin actually exhibits while excluding bank--ZIP cells in which non-entry conveys no information about the bank's choice. Figure~\ref{fig:map_candidate_markets} illustrates this for three large banks (JPMorgan Chase in 2016, Bank of America in 2018, Wells Fargo in 2022): blue ZIPs are those in which the bank operates, gold ZIPs are inside the bank's CBSAs but contain no branch of that bank. The breadth of the gold region relative to the blue makes visible the basic content of the empirical exercise --- the entry model asks what distinguishes opened ZIPs from their gold neighbors within the same CBSAs. All analyses partition the panel into three asset-size buckets ($\$1$B--$\$10$B, $\$10$B--$\$100$B, $\geq \$100$B), motivated by Figure~\ref{fig:openings_by_size_stacked}: the composition of de novo openings has shifted toward the larger buckets over the sample period even as the U.S. branch network contracts overall, and the bank-level edges that drive entry differ across size strata. The resulting panel contains 1{,}166{,}487 bank--ZIP--year observations.

The first channel is local deposit contestability. The branch-closure literature reads contestability as a force that drives banks \emph{out} of rate-sensitive markets, so a credible test on the entry margin requires a measure that varies across the bank--ZIP cells where entry decisions are made. We construct a bank--ZIP--year deposit beta following the within-cycle first-stage approach of \citet{narayanan2025decline} adapted from \citet{DSS2017}. For each of three recent Federal Reserve tightening cycles (2004:Q1--2006:Q1, 2016:Q1--2019:Q1, 2022:Q1--2023:Q1), we measure the bank-level change in the deposit interest-expense ratio between cycle endpoints and regress it within-cycle on deposit-weighted ZIP demographics (age, dividend-receiving share, college share, log family income, log population density, county deposit concentration) and bank balance-sheet structure (log assets, transaction-deposit share, uninsured share, time-deposit share). The within-cycle coefficients are then projected onto every bank--ZIP--year by substituting ZIP demographics and the bank's balance sheet at $t-1$, and divided by the realized cycle rate change to convert into a deposit beta. ZIPs through 2015 inherit the 2004--2006 first-stage; through 2021 the 2016--2019; from 2022 onward the 2022--2023. The resulting measure varies across ZIPs (through local demographics and concentration), across banks within a ZIP (through balance-sheet structure), and across cycles within a bank--ZIP.

The second channel is the bank's cost-side edge against the incumbents it would face in a candidate market. The hypothesis is that a bank with a structurally lower operating cost can monetize a branch in a contestable market in a way an incumbent cannot, even at the same deposit rate; the measurement target is therefore a bank-versus-local-incumbent gap, not a level. We construct two such gaps, each against a self-excluded incumbent benchmark so the focal bank does not contaminate its own reference distribution. The rate-paying advantage subtracts the focal bank's year-$t-1$ efficiency ratio from the deposit-weighted average efficiency ratio of all other banks operating in the ZIP; positive values indicate the focal bank operates more efficiently than the incumbents it would compete with, and the main interactions use a binary indicator for a positive gap. The interest-expense undercut share isolates the funding-side margin: it is the deposit-weighted fraction of self-excluded incumbents in the ZIP whose interest expense per dollar of assets exceeds the focal bank's. A binary indicator for an undercut share above one half is the variant used in the main interactions. Both gaps are winsorized at the one percent tails within year.

The third channel is the bank's existing non-branch lending presence. An entry into a ZIP where the bank already lends turns the branch into a distribution channel for a higher-lifetime-value customer rather than a freestanding deposit-rent vehicle. Mortgage presence is an indicator equal to one if the bank originates any HMDA-reported mortgage in the ZIP in $t-1$; CRA small-business presence is the analogous indicator from the CRA aggregate small-business loan tables. The continuous market-share analogues are constructed for completeness but the binary measures are preferred because the channel of interest is the existence of a non-branch customer relationship, not its scale. Because the opportunity set is restricted to the bank's existing CBSA footprint, both indicators are typically positive in the pooled panel; the channel separates entries from non-entries through size-bucket interactions and through the marginal ZIPs at the edges of the lending footprint, both pursued in Section~\ref{sec:results}.

Every specification conditions on a fixed vector of ZIP-level controls measured at $t-1$. Three flow controls capture local credit and deposit growth (log changes from $t-3$ to $t-1$ in total ZIP deposits, HMDA mortgage originations, and CRA small-business originations). A local-market vector follows: log total ZIP deposits in $t-1$, log average AGI per IRS return, median age, the percent of adults twenty-five and over with a bachelor's degree or higher, log population density, an indicator for ZIPs with no SOD-reporting depository, and an indicator for whether the focal bank has any branch in the ZIP's state. Three existing-footprint controls --- the bank's own deposit share in the ZIP at $t-1$, the log of one plus the number of branches the bank already operates in the ZIP at $t-1$, and the same-state indicator --- absorb the lateral-expansion margin so that the channel coefficients are read off the marginal entry decision rather than the expansion of an already-present footprint.

\subsection{Descriptive Statistics}
\label{sec:data:desc}

Table~\ref{tab:descriptive_stats} reports the sample moments of every variable used in the analysis, separately for each size bucket. Panel A covers banks above one hundred billion dollars in assets and shows the dependent variable, the channel variables, and the full vector of ZIP-level controls; Panels B and C cover the ten-to-hundred-billion and one-to-ten-billion buckets and report only the dependent variable and the channel variables, because the ZIP-level controls have very similar moments across buckets by construction (the candidate-ZIP set is restricted to the bank's CBSA footprint in each case). De novo openings are rare: the opening rate is 0.6 percent in the largest bucket and falls to 0.4 percent and 0.3 percent as we move to mid-sized and small banks. The fall is the candidate-set arithmetic at work --- larger banks span larger CBSA footprints and therefore larger denominators per opening --- and is what makes the size-stratified estimation rather than a pooled regression the natural unit of analysis.

The contestability and cost-side channels each show a distinct cross-bucket pattern. Deposit beta is highest at the largest banks (mean 0.29) and falls monotonically to 0.24 at the smallest, reflecting that large-bank footprints concentrate in dense, demographically high-pass-through markets where the within-cycle first stage assigns more contestability. The rate-paying advantage and the interest-expense undercut share are each near the construction-implied fifty-percent benchmark across all three buckets, but the cross-bucket patterns differ: the rate-paying binary peaks at the mid-bank scale (sixty percent) while the IEA undercut binary rises monotonically as bank size falls, from forty-seven percent in the largest bucket to sixty-two percent in the smallest. Smaller banks are more likely to fund themselves cheaply relative to their local incumbents but less consistently to operate at a lower cost-to-revenue ratio, a pattern the size-stratified channel specifications in Section~\ref{sec:results} are designed to exploit.

The cross-sell channel shows the sharpest size gradient. Mortgage presence is one in eighty percent of bank--ZIP--year cells in the largest bucket, fifty-one percent at mid-sized banks, and thirty-two percent at small banks; CRA small-business presence runs ninety-one, eighty, and fifty-two percent across the same three buckets. The very high cross-sell presence at the largest banks reflects nationally-active mortgage and small-business lending franchises that cover most ZIPs in the bank's CBSA footprint, so cross-sell binaries have little within-bucket separation power for that group; among mid and small banks, the binaries vary across a meaningful share of the candidate ZIPs and therefore carry the lending-presence channel.

The cross-bucket asymmetry across these three channels is the empirical hook for the size-stratified analysis. Variation in the bank's edge against a local market is concentrated in deposit beta at the largest banks (where cost-side and cross-sell binaries are near-saturated or near-fifty-percent), in cross-sell presence at the smaller banks (where deposit beta is more compressed), and in a mix of cost-side and cross-sell measures in between. Section~\ref{sec:results} reads each channel against the bucket in which it has bite, and the interaction specifications surface the channel-by-bucket structure these summary statistics motivate.
